% ============================================================
%  back/appendix-visual-guide.tex
%  Visual Architecture Guide
% ============================================================

\chapter*{Visual Architecture and Diagram Guide}
\addcontentsline{toc}{chapter}{Visual Architecture and Diagram Guide}

\section*{Design principles}

Every diagram in this monograph follows a consistent visual
grammar reflecting the book's central thesis:
every observation is a projection, and projections have kernels.

\subsection*{Partition of diagrams}

Each diagram explicitly separates three zones:
\begin{description}[leftmargin=2cm]
  \item[Hidden layer] The global structure $\mathcal{M}$
    or underlying field $(\Phi, \mathbf{v}, S, \ell)$.
    Rendered in light gray or dashed lines.
  \item[Projection layer] The map $F$ and observation map $\pi$.
    Rendered as arrows with kernel indicated by a shaded band.
  \item[Observable layer] The observable output $\mathcal{O}$.
    Rendered in solid lines.
\end{description}

\subsection*{Phase-space and coarsening diagrams}

\begin{description}[leftmargin=2cm]
  \item[High contrast] Dark background with luminous traces,
    evoking terminal-output clarity.
    Reinforces the algorithmic, foundational character of the theory.
  \item[Boundary encoding] The boundary separating coherent
    ($L \gg D$) and incoherent ($D \gg L$) phases is always
    shown explicitly.
  \item[Kernel shading] The kernel $\ker(\pi \circ F)$ is
    indicated by a shaded region in every projection diagram.
\end{description}

\subsection*{Standard diagram types}

\begin{enumerate}[leftmargin=2em]
  \item \textbf{Coarsening sequence:}
    Four panels showing $t = 0, t_1, t_2, t_\infty$.
    Domain size $\ell(t)$ annotated.
    Defects marked with $\times$.
  \item \textbf{Free-energy landscape:}
    $\mathcal{F}$ vs.\ order parameter, with double-well
    potential $W(\phi)$ and trajectory arrows indicating descent.
  \item \textbf{Projection diagram:}
    $X \xrightarrow{F} L \xrightarrow{\pi} O$ as a commutative
    rectangle with kernel shaded.
  \item \textbf{Sheaf diagram:}
    Open sets $U_1, U_2$ overlapping in $U_{12}$,
    sections $s_1, s_2$ as arrows, mismatch
    $\check{\delta}(s)$ indicated.
  \item \textbf{Cosmic web topology:}
    Simplicial complex $K$ with vertices (clusters),
    edges (filaments), faces (walls), shaded voids.
  \item \textbf{RSVP phase diagram:}
    $Pe$--$\Gamma$ plane with contours of $z_{\mathrm{RSVP}}$.
    Regions labeled: Allen--Cahn regime, Cahn--Hilliard regime,
    advection-dominated, vorticity-stabilised.
  \item \textbf{Redshift path diagram:}
    Null geodesic $\gamma$ through RSVP medium,
    $\mathcal{R}_{\mathrm{eff}}$ colour-coded along path,
    accumulated $\int_\gamma\mathcal{R}\,ds$ plotted below.
\end{enumerate}

\subsection*{The Reconstruction View diagram}

Every chapter containing a \textbf{Reconstruction View} box
is accompanied by a small standardised diagram showing:
\begin{itemize}[leftmargin=2em]
  \item The observable $\mathcal{O}$ (solid box, bottom).
  \item The projection $\pi \circ F$ (downward arrow).
  \item The global structure $\mathcal{M}$ (dashed box, top).
  \item The kernel $\ker(\pi \circ F)$ (shaded region beside arrow).
  \item The reconstruction $G$ (upward dashed arrow).
\end{itemize}

This diagram repeats in every chapter where the Reconstruction
View box appears, with only the labels changing.
By Chapter~41, the reader should be able to sketch it from memory.
